\section{工件索引}
\label{app:artifacts}

\begin{table}[H]
\centering
\footnotesize
\setlength{\tabcolsep}{4pt}
\renewcommand{\arraystretch}{0.92}
\caption{支撑本文每个数字的工件。}
\label{tab:artifacts}
\begin{tabular}{ll}
\toprule
项目 & 路径 \\
\midrule
响应性审计 & \url{scripts/78_responsiveness_audit.py} \\
选择输入适配器 & \url{scripts/79_export_selective_inputs.py} \\
绘图脚本 & \url{scripts/80_make_figures.py} \\
敏感性分析 & \url{scripts/81_sensitivity_analyses.py} \\
图像侧可检测性 & \url{scripts/82_image_side_baseline.py} \\
质量语言审计 & \url{scripts/83_quality_language_audit.py} \\
PadChest 疾病阳性子集 & \url{scripts/84_padchest_abnormal_subset.py} \\
工作点区间 & \url{scripts/85_operating_point_ci.py} \\
患病率分层 & \url{scripts/87_label_hierarchy_prevalence.py} \\
逐标签 CI / BH-FDR & \url{scripts/88_per_label_ci_fdr.py} \\
固定阈值事后分析 & \url{scripts/89_fixed_threshold_posthoc.py} \\
覆盖普查 & \url{scripts/90_coverage_census.py} \\
清单锁定 & \url{scripts/91_manifest_lock.py} \\
round-2 汇总 & \url{scripts/96_analyze_round2.py} \\
配对阳性率 & \url{scripts/102_paired_positive_rate.py} \\
DenseNet 患者自助法 & \url{scripts/103_densenet_bootstrap.py} \\
表格一致性核查 & \url{scripts/104_table_consistency_check.py} \\
输出清单 & \url{scripts/105_outputs_manifest.py} \\
干净 AUROC 区间 & \url{scripts/106_clean_auroc_ci.py} \\
PadChest 标签层级 & \url{scripts/107_padchest_label_hierarchy.py} \\
固定规则重跑 & \url{scripts/108_fixed_rule_rerun.py} \\
CheXbert 交叉核验 & \url{scripts/109_chexbert_openi.py} \\
贪婪确定性 & \url{scripts/110_compare_determinism.py} \\
DenseNet 19 格网格 & \url{scripts/111_densenet_grid.py} \\
跨家族 schema & \url{scripts/112_analyze_xfamily_schema.py} \\
DenseNet 条件指标 & \url{scripts/113_densenet_conditions.py} \\
标注器一致性 & \url{scripts/114_report_labeller_agreement.py} \\
LLaVA-Med NIH 子集 & \url{scripts/115_llava_nih.py} \\
跨家族改写 & \url{scripts/116_paraphrase_xfamily.py} \\
确定性输出 & \url{outputs/determinism/medgemma_4b_it/determinism.csv} \\
固定规则输出 & \url{outputs/fixed_rule/} \\
LLaVA-Med NIH 运行 & \url{outputs/llava_nih/} \\
敏感性结果表 & \url{outputs/analysis_sensitivity/} \\
响应性结果表 & \url{outputs/analysis_responsiveness/} \\
LLaVA-Med 响应性 & \url{outputs/analysis_responsiveness_llava/} \\
契约覆盖 & \url{outputs/analysis_selective/export_manifest.json} \\
地理迁移汇总 & \url{outputs/geographic/<model>/} \\
选择性指标 & \url{outputs/selective_v2/}, \url{outputs/trust_v1/} \\
风险--覆盖带 & \url{outputs/selective_v2/<model>/risk_coverage_bands.csv} \\
跨参数挖掘 & \url{outputs/analysis_nih_mine/joint_*.csv} \\
运行状态（38 阶段） & \url{outputs/state/run_all_layers_state.json} \\
预处理契约 & \url{docs/55_preprocessing_contract.md} \\
\bottomrule
\end{tabular}
\end{table}

\section{协议表}
\label{app:protocol}

\begin{table}[H]
\centering
\small
\caption{NIH shift 网格上的解析与契约覆盖（每模型 28,614 次调用）。``全部确定''是解析成功回答中八个状态全为阳性/阴性的比例。完整 41 阶段普查（交付的 39 个阶段加两个 round-2 推理阶段）见 \url{outputs/analysis_sensitivity/coverage_census.csv}。}
\label{tab:coverage}
\begin{tabular}{lrrr}
\toprule
模型 & 解析率 & 全部确定 & 错误行 \\
\midrule
Qwen2.5-VL & 1.000 & \textbf{1.000} & 0 \\
InternVL3.5 & 1.000 & 0.956 & 0 \\
MedGemma-1.5 & 0.998 & 0.824 & 0 \\
Qwen3-VL & 1.000 & 0.795 & 0 \\
MedGemma-4B & 1.000 & 0.605 & 0 \\
CheXagent & 0.985 & \textbf{0.096} & 0 \\
\bottomrule
\end{tabular}
\end{table}

\begin{table}[H]
\centering
\small
\caption{同一批 947 例 Open-i 研究在两种标注器下的结果（均为槽位级：八个标签 $\times$ 947 次判定）。绝对计数不可跨标注器比较，但排序可以。幻觉排序方向保持一致（Spearman 0.71，$n=6$，$p=0.11$），漏报排序不一致（$-0.49$，$p=0.33$）。}
\label{tab:chexbert}
\begin{tabular}{lrrrr}
\toprule
模型 & 规则漏报 & 规则幻觉 & CheXbert 漏报 & CheXbert 幻觉 \\
\midrule
MedGemma-4B & 196 & \textbf{3,021} & 118 & 609 \\
MedGemma-1.5 & 571 & 268 & 103 & 583 \\
InternVL3.5 & 274 & 2,198 & 187 & 870 \\
Qwen3-VL & \textbf{181} & \textbf{3,508} & 363 & 786 \\
Qwen2.5-VL & 302 & 2,128 & 231 & 660 \\
CheXagent & \textbf{1,198} & \textbf{113} & 165 & 532 \\
\bottomrule
\end{tabular}
\end{table}

\section{有效输入几何}
\label{app:geometry}

\begin{table}[H]
\centering
\small
\caption{最大 PadChest 图像（15.37\,MP）上的有效输入几何。同一张胸片变成六种不同输入；这是测量声明，不是中性设置。}
\label{tab:geometry}
\begin{tabular}{llrl}
\toprule
模型 & 有效输入 & 图像 token 数 & 缩放机制 \\
\midrule
Qwen3-VL-8B & 15.24\,MP & 14,884 & 契约提示 \\
Qwen2.5-VL-7B & 12.85\,MP（上限） & 16,384 & 检查点配置 \\
MedGemma-4B/1.5 & 0.80\,MP & 256 & 固定处理器 \\
InternVL3.5-8B & 0.20\,MP & 256 & $1\times1$ 切片网格 \\
CheXagent-8B & 0.20\,MP & 0（交叉注意力） & 固定处理器 \\
LLaVA-Med-7B & 0.11\,MP & 576 & 固定 CLIP \\
\bottomrule
\end{tabular}
\end{table}

\section{敏感性分析}
\label{app:risk}

所有数值由脚本 81--85（表~\ref{tab:artifacts}）从冻结预测文件与原始图像生成。

\begin{table}[H]
\centering
\small
\caption{众数占比阈值扫描下的坍缩标记（唯一回答数上限固定为 5）。跨数据集判定在 0.40 以上不变；只有在最宽松的 0.30 切点，PadChest 的第六个模型才加入，且只有域内与散文单元在边际上移动。}
\label{tab:sweep}
\begin{tabular}{lrrrrrr}
\toprule
数据层 & 0.30 & 0.40 & 0.50 & 0.60 & 0.70 & 0.80 \\
\midrule
NIH（共 6） & 3 & 3 & 3 & 2 & 1 & 1 \\
PadChest（共 6） & 6 & 5 & 5 & 5 & 5 & 5 \\
Open-i（共 6） & 2 & 1 & 1 & 1 & 0 & 0 \\
\bottomrule
\end{tabular}
\end{table}

\begin{table}[H]
\centering
\footnotesize
\setlength{\tabcolsep}{3pt}
\caption{两个跨域统计量的敏感性。共同标签 $\Delta$AUROC 排除肺炎（在 PadChest 中结构性缺失）。零模型在随机真实图像对与参考长度匹配真实图像对下重算。CheXagent 只有一个可估计的共同标签（记为``(1)''）。最后一列是两种零模型下 $\Delta_{\text{null}}$ 的范围；正值表示空白干预比随机真实图像离真实图像更远。}
\label{tab:sens}
\begin{tabular}{lrrrrrr}
\toprule
模型 & 平均 $\Delta$（7 标签） & 最差标签 & $J$(真实,空白) & 零模型（随机） & 零模型（长度） & $\Delta_{\text{null}}$ 范围 \\
\midrule
MedGemma-4B & $-0.107$ & $-0.246$ & 0.894 & 0.824 & 0.844 & $-0.070$ 至 $-0.050$ \\
MedGemma-1.5 & $-0.079$ & $-0.161$ & 0.736 & 0.648 & 0.654 & $-0.087$ 至 $-0.081$ \\
InternVL3.5 & $-0.134$ & $-0.274$ & 0.562 & 0.673 & 0.677 & $+0.112$ 至 $+0.116$ \\
Qwen3-VL & $-0.059$ & $-0.185$ & 0.091 & 0.740 & 0.734 & $+0.643$ 至 $+0.648$ \\
CheXagent & $-0.004$ (1) & $-0.004$ & 0.277 & 0.247 & 0.252 & $-0.030$ 至 $-0.024$ \\
Qwen2.5-VL & $-0.002$ & $-0.008$ & 0.576 & 0.770 & 0.767 & $+0.191$ 至 $+0.193$ \\
\bottomrule
\end{tabular}
\end{table}

\begin{table}[H]
\centering
\footnotesize
\setlength{\tabcolsep}{3pt}
\caption{标签层级部分受患病率混杂。A 面板：每个模型的保持率与 NIH 患病率 Spearman，以及高/低患病率标签分层的平均保持率。B 面板：秩残差化患病率前后的两两一致性。原始一致性为正，但偏一致性在不同模型对间变号，因此我们不主张患病率无关的病灶规律。}
\label{tab:hierarchy}
\begin{tabular}{lrrr}
\toprule
\multicolumn{4}{l}{\emph{A 面板：保持率与患病率}} \\
模型 & $\rho$(保持率, 患病率) & 平均保持率（高患病率） & 平均保持率（低患病率） \\
\midrule
InternVL3.5 & 0.571 & 0.935 & 0.764 \\
Qwen3-VL & 0.429 & 0.718 & 0.629 \\
MedGemma-1.5 & 0.714 & 0.873 & 0.701 \\
MedGemma-4B & 0.464 & 0.735 & 0.514 \\
\midrule
\multicolumn{4}{l}{\emph{B 面板：模型对一致性}} \\
模型对 & 原始 $\rho$ & 偏 $\rho$ & 精确 $p$ \\
\midrule
InternVL / Qwen3 & 0.771 & 0.765 & 0.658 \\
InternVL / MedGemma-1.5 & 0.857 & $-0.654$ & 0.012 \\
InternVL / MedGemma-4B & 0.750 & $-0.685$ & 0.200 \\
Qwen3 / MedGemma-1.5 & 0.657 & 0.525 & 0.058 \\
Qwen3 / MedGemma-4B & 0.657 & $-0.691$ & 0.242 \\
MedGemma-1.5 / MedGemma-4B & 0.750 & 0.847 & 0.200 \\
\bottomrule
\end{tabular}
\end{table}

\begin{table}[H]
\centering
\footnotesize
\setlength{\tabcolsep}{3pt}
\caption{逐标签干净 $\rightarrow$ 重度召回变化，附患者聚类自助法区间（1,000 次重采样）与每个模型内七个标签的 Benjamini--Hochberg $q$。重度单元合并六种退化类型。心脏增大与接近零召回的气胸是不显著例外；28 个单元中 22 个 $q<0.05$。完整 CSV：\url{outputs/analysis_sensitivity/per_label_bootstrap_fdr.csv}。}
\label{tab:perlabel}
\begin{tabular}{llrrr}
\toprule
模型 & 标签 & $\Delta$ 召回（pp） & 95\% CI & $q_{\mathrm{BH}}$ \\
\midrule
InternVL3.5 & 肺不张 & $-4.23$ & $[-5.71,-2.65]$ & $<0.001$ \\
 & 心脏增大 & $+0.19$ & $[-1.63,+2.18]$ & 0.878 \\
 & 实变 & $-3.04$ & $[-4.85,-1.18]$ & $<0.001$ \\
 & 水肿 & $-8.43$ & $[-12.04,-4.71]$ & $<0.001$ \\
 & 胸腔积液 & $-2.40$ & $[-3.83,-1.16]$ & $<0.001$ \\
 & 肺炎 & $-5.91$ & $[-8.71,-3.11]$ & $<0.001$ \\
 & 气胸 & $-0.08$ & $[-0.81,+0.29]$ & 0.878 \\
\midrule
Qwen3-VL & 肺不张 & $-4.08$ & $[-5.73,-2.42]$ & $<0.001$ \\
 & 心脏增大 & $-2.33$ & $[-6.71,+2.21]$ & 0.401 \\
 & 实变 & $-5.54$ & $[-7.98,-3.29]$ & $<0.001$ \\
 & 水肿 & $-11.98$ & $[-17.13,-7.27]$ & $<0.001$ \\
 & 胸腔积液 & $-4.02$ & $[-5.64,-2.38]$ & $<0.001$ \\
 & 肺炎 & $-10.03$ & $[-14.26,-6.77]$ & $<0.001$ \\
 & 气胸 & $0.00$ & $[0.00,0.00]$ & 1.000 \\
\midrule
MedGemma-1.5 & 肺不张 & $-4.23$ & $[-5.08,-3.40]$ & $<0.001$ \\
 & 心脏增大 & $-5.56$ & $[-7.55,-3.67]$ & $<0.001$ \\
 & 实变 & $-5.05$ & $[-7.67,-2.43]$ & $<0.001$ \\
 & 水肿 & $-17.32$ & $[-20.80,-13.90]$ & $<0.001$ \\
 & 胸腔积液 & $-13.57$ & $[-15.55,-11.71]$ & $<0.001$ \\
 & 肺炎 & $-13.08$ & $[-17.58,-8.57]$ & $<0.001$ \\
 & 气胸 & $-2.83$ & $[-4.78,-1.12]$ & $<0.001$ \\
\midrule
MedGemma-4B & 肺不张 & $-3.49$ & $[-4.40,-2.68]$ & $<0.001$ \\
 & 心脏增大 & $-9.44$ & $[-11.54,-7.49]$ & $<0.001$ \\
 & 实变 & $-11.22$ & $[-14.09,-8.44]$ & $<0.001$ \\
 & 水肿 & $-31.43$ & $[-34.90,-28.06]$ & $<0.001$ \\
 & 胸腔积液 & $-14.53$ & $[-16.85,-12.33]$ & $<0.001$ \\
 & 肺炎 & $-21.68$ & $[-26.51,-17.48]$ & $<0.001$ \\
 & 气胸 & $-0.13$ & $[-0.54,+0.11]$ & 0.746 \\
\bottomrule
\end{tabular}
\end{table}

\begin{table}[H]
\centering
\small
\caption{三个工作点单元的事后固定阈值检查。``状态''使用模型输出的阳性/阴性状态；``$p\ge0.5$''对口头概率施加固定切点；``干净匹配''在干净图像上选一个切点，使其阳性数与状态规则相同。三种规则都在 0.11\,\pp{} 内复现召回率与阳性预测率变化，因此该移动不是离散状态规则改变造成的伪影。表~\ref{tab:fixedrerun} 给出固定规则重跑，用于分离阈值移动与分数分布移动。}
\label{tab:fixedthr}
\begin{tabular}{llrrr}
\toprule
单元 & 规则 & $\Delta$ 召回 & $\Delta$ 精确率 & $\Delta$ 阳性率 \\
\midrule
InternVL 对比度 & 状态 & $+4.69$ & $-1.38$ & $+4.06$ \\
 & $p\ge0.5$ & $+4.80$ & $-1.35$ & $+4.13$ \\
 & 干净匹配 & $+4.69$ & $-1.38$ & $+4.06$ \\
\midrule
InternVL 亮度 & 状态 & $+2.52$ & $-0.64$ & $+2.09$ \\
 & $p\ge0.5$ & $+2.47$ & $-0.78$ & $+2.15$ \\
 & 干净匹配 & $+2.52$ & $-0.59$ & $+2.06$ \\
\midrule
Qwen3-VL 对比度 & 状态 & $+3.44$ & $+0.19$ & $+2.51$ \\
 & $p\ge0.5$ & $+3.36$ & $+0.38$ & $+2.40$ \\
 & 干净匹配 & $+3.44$ & $+0.25$ & $+2.49$ \\
\bottomrule
\end{tabular}
\end{table}

\begin{table}[H]
\centering
\small
\caption{图像侧可检测性。单一无参考统计量把 1,506 张干净的域内评估图像与其退化版本分开；AUC 已做方向调整，1.000 表示完全分离。每个条件只列出 AUC 最高的统计量。退化在图像层面一目了然；静默失效结果只涉及模型侧通道。}
\label{tab:imageside}
\begin{tabular}{llrrrr}
\toprule
条件 & 统计量 & AUC & 方向 & 干净中位数 & 重度中位数 \\
\midrule
噪声（重度） & Laplacian 方差 & \textbf{1.000} & 退化后更高 & 0.0011 & 0.0844 \\
模糊（重度） & Laplacian 方差 & \textbf{1.000} & 退化后更低 & 0.0011 & 0.000024 \\
对比度（重度） & 强度熵 & 0.966 & 退化后更低 & 7.451 & 6.546 \\
对比度（重度） & 强度标准差 & 0.900 & 退化后更高 & 0.231 & 0.320 \\
\bottomrule
\end{tabular}
\end{table}

\begin{table}[H]
\centering
\small
\caption{病例级转移，InternVL3.5，干净 $\rightarrow$ 对比度/重度。``获得''是干净阴性图像在退化下被判为阳性；``新假阳性''是干净时正确的阴性被判为阳性。}
\label{tab:flip}
\begin{tabular}{lrrrrrr}
\toprule
发现 & 保留 TP & 丢失 TP & 获得 TP & 纠正 FP & 新 FP & 稳定 TN \\
\midrule
肺不张 & 535 & 10 & 43 & 5 & 36 & 256 \\
心脏增大 & 324 & 3 & 30 & 28 & 59 & 495 \\
实变 & 99 & 1 & 50 & 8 & 69 & 770 \\
水肿 & 135 & 0 & 23 & 9 & 40 & 993 \\
胸腔积液 & 592 & 19 & 34 & 5 & 21 & 311 \\
肺炎 & 57 & 1 & 22 & 8 & 98 & 1,034 \\
气胸 & 0 & 1 & 2 & 1 & 0 & 1,108 \\
无发现 & 83 & 6 & 1 & 53 & 11 & 1,128 \\
\bottomrule
\end{tabular}
\end{table}

\begin{table}[H]
\centering
\small
\caption{按标签的病例构成（阳性率，\%）。NIH 自然值为官方 ChestX-ray14 在全部 112,120 张图像上的标注；1,506 张评估集由分层采样器做疾病富集。PadChest 为标签映射后的原生阳性率；肺炎结构性缺失。因此精确率与阳性预测率只在层内比较。}
\label{tab:prevalence}
\begin{tabular}{lrrr}
\toprule
发现 & NIH 评估集 & NIH 自然值 & PadChest 原生值 \\
\midrule
肺不张 & 50.7 & 10.3 & 2.8 \\
心脏增大 & 29.0 & 2.5 & 5.8 \\
实变 & 36.0 & 4.2 & 5.5 \\
水肿 & 17.5 & 2.1 & 0.7 \\
胸腔积液 & 66.3 & 11.9 & 2.5 \\
肺炎 & 13.7 & 1.3 & 0.0 \\
气胸 & 26.4 & 4.7 & 0.3 \\
无发现 & 8.0 & 53.8 & 16.8 \\
\bottomrule
\end{tabular}
\end{table}

\begin{table}[t]
\centering
\footnotesize
\setlength{\tabcolsep}{4pt}
\caption{PadChest 原生分辨率下的逐标签 AUROC（$n=1{,}500$；CheXagent 为 1,499；1,000 次患者聚类重采样）。五个无响应模型最多给出 20 种不同回答，其 AUROC 在 0.500 的 $\pm0.04$ 以内，不构成有意义的标签排序。只有 MedGemma-1.5 可估计：心脏增大 0.541（95\% CI 0.505--0.584）、肺不张 0.531（0.464--0.605）；其余四个可估计标签为 0.489--0.490。肺炎结构性缺失。}
\label{tab:padlabel}
\begin{tabular}{lrrrrrr}
\toprule
发现 & MG-1.5 & MG-4B & InternVL & Qwen3 & Qwen2.5 & CheXagent \\
\midrule
肺不张 & 0.531 & 0.505 & 0.500 & 0.499 & 0.500 & 0.500 \\
心脏增大 & \textbf{0.541} & 0.492 & 0.500 & 0.500 & 0.500 & 0.500 \\
实变 & 0.490 & 0.503 & 0.500 & 0.500 & 0.500 & 0.500 \\
水肿 & 0.489 & 0.536 & 0.500 & 0.500 & 0.500 & 0.500 \\
胸腔积液 & 0.489 & 0.516 & 0.500 & 0.500 & 0.500 & 0.500 \\
肺炎 & --- & --- & --- & --- & --- & --- \\
气胸 & 0.489 & 0.492 & 0.500 & 0.500 & 0.500 & 0.500 \\
\bottomrule
\end{tabular}
\end{table}

\begin{table}[t]
\centering
\small
\caption{DenseNet-121（表~\ref{tab:baseline} 的检查点）在 19 格网格上以固定 0.5 规则运行。对比度 severe 提高召回与阳性率而宏 AUROC 保持平稳，说明工作点移动并不需要语言模型；噪声与分辨率使两者同时下降。PadChest 宏 AUROC 为 0.515，因此跨数据集坍缩也不是 VLM 特有的。}
\label{tab:densenetgrid}
\begin{tabular}{lrrr}
\toprule
格子 & 宏 AUROC & 合并召回 & 阳性率 \\
\midrule
clean & 0.7329 & 0.5673 & 0.3599 \\
对比度 severe & 0.7292 & \textbf{0.5878} & \textbf{0.3775} \\
模糊 severe & 0.7296 & 0.5655 & 0.3585 \\
亮度 severe & 0.7109 & 0.5315 & 0.3602 \\
噪声 severe & 0.6824 & 0.4766 & 0.3208 \\
分辨率 severe & 0.7107 & 0.5230 & 0.3548 \\
PadChest（原生） & 0.5153 & 0.3232 & --- \\
\bottomrule
\end{tabular}
\end{table}

\begin{table}[H]
\centering
\small
\caption{200 例配对研究上的固定规则重跑（InternVL3.5）。模型被重新提示为只输出概率，阳性/阴性状态由我们在 0.5 处赋予。由于该运行的决策规则不可能移动，两种退化下召回与阳性率的上升来自分数分布移动，而不是离散状态改变。Qwen3-VL 在同一``只输出概率''契约下对每个标签都返回 0.0（非零率 1.0\%，0.5 切点阳性率 $\le$ 1.0\%），因此该单元退化、被排除。}
\label{tab:fixedrerun}
\begin{tabular}{lrrrrr}
\toprule
单元 & $n$ & 召回（干净） & 召回（重度） & 阳性率（干净） & 阳性率（重度） \\
\midrule
对比度 severe & 200 & 0.510 & \textbf{0.575} & 0.404 & \textbf{0.449} \\
亮度 severe & 200 & 0.510 & \textbf{0.584} & 0.404 & \textbf{0.457} \\
\bottomrule
\end{tabular}
\end{table}

\section{响应性审计的逐单元证据}
\label{app:cell}

审计对每个（数据层，模型）报告：真实图像上的输出支撑大小、带 Wilson 区间的众数回答占比、结构化层中不同\emph{状态向量}的数量、空白模板集中度，以及两个零模型差距。完整表见 \url{outputs/analysis_responsiveness/responsiveness_by_cell.csv}；表~\ref{tab:audit} 是其精简形式。坍缩规则（众数占比 $\ge 0.50$ 或唯一回答数 $\le 5$）在脚本头部声明，仅作描述：没有任何一行被它删除或重新缩放。

\section{确证性实验表}
\label{app:round2}

\subsection*{第二个医学模型上的 schema 扫描}

\begin{table}[t]
\centering
\footnotesize
\setlength{\tabcolsep}{3.5pt}
\caption{schema 扫描，MedGemma-1.5（每格 $n=1{,}506$，共 15,060 次调用）。``空白召回''为均匀灰色图像上七个病理标签的合并召回；``违约''为至少一个标签缺少明确状态或概率的回答比例。}
\label{tab:r2schema}
\begin{tabular}{lrrrrr}
\toprule
schema & 真实召回 & 空白召回 & 空白 \texttt{no\_finding} 阳性 & 空白违约 & 真实违约 \\
\midrule
v0 简单 & 0.551 & \textbf{0.000} & 1.000 & 0.000 & 0.059 \\
v1 逐标签二分类 & 0.493 & \textbf{0.000} & 1.000 & 0.000 & 0.221 \\
v2 few-shot & 0.550 & \textbf{0.000} & 0.001 & 0.999 & 0.250 \\
v3 先报告后标注 & 0.552 & \textbf{0.000} & 1.000 & 0.000 & 0.152 \\
v4 正常性施压 & \textbf{0.614} & \textbf{0.000} & 0.000 & 1.000 & 0.531 \\
\bottomrule
\end{tabular}
\end{table}

\subsection*{改写与 v0 基线的一致性}

\begin{table}[H]
\centering
\footnotesize
\setlength{\tabcolsep}{3pt}
\caption{改写对照（每格 $n=1{,}506$）。MedGemma-4B 六格、6,024 次调用；Qwen3-VL 同批研究四格。一致率与 $\kappa$ 按标签与同模型 v0 比较后对八个标签取平均；单一状态时 $\kappa$ 无定义。}
\label{tab:r2para}
\begin{tabular}{llrrr}
\toprule
模型 & 图像 / 改写 & 一致率 & $\kappa$ & \texttt{no\_finding} 一致率 \\
\midrule
MedGemma-4B & 真实 / p1 简短 & 0.909 & 0.624 & 0.999 \\
 & 真实 / p2 第二读者 & 0.914 & 0.642 & 0.946 \\
 & 空白 / p1 简短 & 0.875 & 0.000 & \textbf{0.000} \\
 & 空白 / p2 第二读者 & 1.000 & 0.000 & \textbf{1.000} \\
\midrule
Qwen3-VL & 真实 / p1 简短 & 0.931 & 0.723 & 0.926 \\
 & 真实 / p2 第二读者 & 0.903 & 0.701 & 0.909 \\
 & 空白 / p1 简短 & 1.000 & --- & \textbf{1.000} \\
 & 空白 / p2 第二读者 & 1.000 & --- & \textbf{1.000} \\
\bottomrule
\end{tabular}
\end{table}

\subsection*{第二个推理模型的终止预算}

\begin{table}[H]
\centering
\small
\caption{Qwen3-VL-8B-Thinking 的预算升级（10,542 个格子）。前三行为全量运行；第四行并入仍失败格子的 40 格随机抽样。``中位长度''指该预算下仍无可用答案的回复；它随上限上升，说明这些回复会撑满给它的任何预算。计数为格子，不是模型调用。}
\label{tab:r2budget}
\begin{tabular}{lrrrr}
\toprule
预算（token） & 可用格子 & 覆盖率 & 仍失败 & 中位长度（千字符） \\
\midrule
2,560 & 6,067 & 57.6\% & 4,475 & 9.8 \\
3,584 & 8,738 & 82.9\% & 1,804 & 13.5 \\
5,120 & 9,381 & 89.0\% & 1,161 & 19.3 \\
5,120 + 8,192 抽样 & 9,401 & 89.2\% & 1,141 & 28.1 \\
\bottomrule
\end{tabular}
\end{table}

\subsection*{失败归档与残余格子的对账}

\begin{table}[H]
\centering
\small
\caption{失败归档的构成。行是\emph{尝试}，因此一个格子在更大预算下重试就会在每个预算上各出现一次；最终过滤后的格子级残余是 1,141，而不是归档行数。}
\label{tab:r2recon}
\begin{tabular}{llrr}
\toprule
轮次标签 & 预算（token） & 归档行数 & 不同格子数 \\
\midrule
\texttt{pre\_5120} & 2,560 / 3,584 & 3,085 & 2,491 \\
\texttt{post\_5120} & 5,120 & 1,161 & 1,161 \\
\texttt{cap8192} & 8,192 & 37 & 20 \\
\midrule
合计 & --- & \textbf{4,283} & 2,761 \\
\bottomrule
\end{tabular}
\end{table}

这一区分对报告很重要：``1,141 个格子在 5,120 下从不终止''是格子级陈述，而``4,283 条归档失败''跨四个预算计尝试数。若把二者等同，会把残余高估 3.8 倍。

\subsection*{配对工作点（控制覆盖率）}

\begin{table}[H]
\centering
\small
\caption{Qwen3-VL-8B-Thinking 在配对子集上的阳性率：在 clean 格与 severe 格都能给出可用答案的研究。配对与非配对的移动一致，因此工作点移动不是由差异化覆盖率造成的。}
\label{tab:r2posrate}
\begin{tabular}{lrrrr}
\toprule
格子 & 配对 $n$ & clean & severe & $\Delta$ \\
\midrule
blur & 1,302 & 0.249 & 0.192 & $-5.7$\,\pp{} \\
brightness & 1,178 & 0.255 & 0.215 & $-4.0$\,\pp{} \\
contrast & 1,299 & 0.254 & 0.278 & $+2.4$\,\pp{} \\
JPEG & 1,290 & 0.254 & 0.243 & $-1.0$\,\pp{} \\
noise & 1,238 & 0.249 & 0.172 & $-7.7$\,\pp{} \\
resolution & 1,132 & 0.250 & 0.128 & $\mathbf{-12.2}$\,\pp{} \\
\bottomrule
\end{tabular}
\end{table}

\subsection*{第二个推理模型的逐标签损伤}

\begin{table}[H]
\centering
\small
\caption{Qwen3-VL-8B-Thinking 的 clean $\rightarrow$ 合并 severe 逐标签召回变化，在 clean 与全部六个 severe 格都有答案的 1,394 例研究（695 名患者）上配对；逐格配对 $n$ 从 1,130（resolution）到 1,297（blur），逐格估计见表~\ref{tab:r2think}。区间为患者聚类自助法（1,000 次重采样，95\%）；$q$ 为七个标签上的 Benjamini--Hochberg 值。}
\label{tab:r2label}
\begin{tabular}{lrrrrr}
\toprule
标签 & clean & severe & $\Delta$ & 95\% CI & $q$ \\
\midrule
肺不张 & 0.089 & 0.067 & $-2.2$\,\pp{} & $[-4.4,-0.2]$ & 0.053 \\
心影增大 & 0.592 & 0.532 & $-6.0$\,\pp{} & $[-9.9,-1.9]$ & 0.014 \\
实变 & 0.314 & 0.288 & $-2.6$\,\pp{} & $[-6.5,+1.1]$ & 0.166 \\
肺水肿 & 0.117 & 0.075 & $-4.2$\,\pp{} & $[-8.5,-0.2]$ & 0.053 \\
胸腔积液 & 0.482 & 0.394 & $\mathbf{-8.8}$\,\pp{} & $[-11.9,-5.8]$ & $<0.001$ \\
肺炎 & 0.451 & 0.341 & $\mathbf{-11.0}$\,\pp{} & $[-16.4,-5.5]$ & $<0.001$ \\
气胸 & 0.003 & 0.001 & $-0.1$\,\pp{} & $[-0.8,+0.3]$ & 0.820 \\
\bottomrule
\end{tabular}
\end{table}
